\section{Evidencia visual y cronología de la investigación}

\subsection{Slide 3: Pipeline del modelo de Q-Learning}

Esta página muestra el pipeline de Q-Learning desde los datos (from data) hasta el despliegue (deployment); cada etapa requiere un guardrail (replay hygiene, target network, threshold monitor).

\begin{figure}[H]
\centering
\includegraphics[width=0.92\textwidth,page=3]{06_cs224r_qlearning_2025.pdf}
\caption{Slide: Q-Learning pipeline and guardrail checks}
\end{figure}
\newpage

\subsection{Slide 4: Temporal difference y replay dynamics}

La Slide 4 relaciona el TD error con el replay sampling y enfatiza que el muestreo de `high-error` contribuye al prioritized buffer, pero que también requiere revisión humana.

\begin{figure}[H]
\centering
\includegraphics[width=0.92\textwidth,page=4]{06_cs224r_qlearning_2025.pdf}
\caption{Slide: TD error monitoring y prioritized sampling}
\end{figure}
\newpage

\subsection{Slide 5: Monitoreo del deployment}

Durante el despliegue se necesita un monitoreo multicanal de la reward distribution, la episode success rate y la latency; la Slide 5 muestra el telemetry stack.

\begin{figure}[H]
\centering
\includegraphics[width=0.92\textwidth,page=5]{06_cs224r_qlearning_2025.pdf}
\caption{Slide: deployment telemetry and alert stack}
\end{figure}
\newpage

\subsection{Slide 6: Action timeline and evidence matrix}

La Slide 6 incluye la evidence matrix y explica que cada failure tiene su action timeline correspondiente, tal como muestra la table que agregamos en este documento.

\begin{figure}[H]
\centering
\includegraphics[width=0.92\textwidth,page=6]{06_cs224r_qlearning_2025.pdf}
\caption{Slide: Evidence matrix and action timeline}
\end{figure}
\newpage

\subsection{Slide 7: Research directions}

Esta página señala las future directions: multi-step Q-Learning, safety evaluation, sim-to-real.

\begin{figure}[H]
\centering
\includegraphics[width=0.92\textwidth,page=7]{06_cs224r_qlearning_2025.pdf}
\caption{Slide: Research directions for safe Q-Learning}
\end{figure}
\newpage

\subsection{Slide 8: Time/Loss visualization}

La Slide 8 muestra el descenso simultáneo de la reward curve y del TD error, lo que ofrece una pista visual para cerrar el documento.

\begin{figure}[H]
\centering
\includegraphics[width=0.92\textwidth,page=8]{06_cs224r_qlearning_2025.pdf}
\caption{Slide: Reward \& TD error visualization}
\end{figure}
\newpage

